\section{Experiments and In-Depth Diagnostic Analysis}
\label{sec:experiments}

In this section, we present a rigorous empirical evaluation of \textbf{FluidMerge} alongside three state-of-the-art baselines: SyMerge \cite{jung2025symerge}, AdaMerging \cite{yang2024adamerging}, and Task Surgery \cite{zhang2024task}. We do not merely present metrics; we systematically evaluate performance under standardized protocols and diagnose the deep representation gaps and physical failure modes of post-hoc test-time model adaptation.

\subsection{Experimental Setup}

\paragraph{Architecture and Checkpoints}
We evaluate all methods using a Vision Transformer backbone (\textbf{ViT-B-32}) \cite{dosovitskiy2020image, radford2021clip} containing approximately 86 million parameters. The evaluation features eight diverse image classification datasets: SUN397, Stanford Cars, RESISC45, EuroSAT, SVHN, GTSRB, MNIST, and DTD. Critically, we utilize pre-trained task-specific classification heads (\texttt{head\_<dataset>.pt}) that were fine-tuned on top of fully fine-tuned expert checkpoints. 

\paragraph{Method-Specific Configuration}
\begin{itemize}
    \item \textbf{FluidMerge (Ours):} We run the continuous-time physical simulation for $N=100$ steps with a step size of $\Delta t = 0.1$ ($T=10.0$). The viscosity coefficient is set to $\nu = 0.001$. The classifier heads are jointly trained via an Adam optimizer with a learning rate of $10^{-2}$. Because our task-aligned data-routing cleanly feeds each expert teacher its native, in-distribution stream, cross-task out-of-distribution noise is completely prevented by design. Thus, the confidence-based entropy filtering threshold is set to $\tau = \infty$ (disabled) for our primary classification experiments, as filtering is not required when OOD inputs are bypassed structurally.
    \item \textbf{SyMerge Baseline:} Evaluated over 500 epochs, adapting task-specific layers using teacher-guided self-labeling.
    \item \textbf{AdaMerging Baseline:} Evaluated over 500 epochs, learning layer-wise coefficients via entropy minimization.
    \item \textbf{Task Surgery Baseline:} Evaluated over 1000 epochs, performing L1-based feature alignment.
\end{itemize}
To ensure strict experimental standardization and exact reproducibility, all methods are optimized under identical test-time adaptation (TTA) batch settings. Specifically, for each of the eight evaluation datasets, we randomly sample $1000$ unlabeled images from the respective validation or test splits. During the continuous-time physical simulation and baseline TTA adaptation loops, these unlabeled images are fed to the model using a standard mini-batch size of $32$ samples.

\subsection{Main Comparative Evaluation: The Synergy-Refinement Protocol}
\label{subsec:synergy_refinement}

To ensure a mathematically rigorous, fair comparison under identical, high-performing initialization conditions, we evaluate all test-time adaptation baselines—including SyMerge \cite{jung2025symerge} and AdaMerging \cite{yang2024adamerging}—under our proposed \textbf{Synergy-Refinement Protocol}. In this protocol, all dynamic continuous-time simulations and adaptation baselines are initialized starting from the \textbf{Expert-Weighted Initial Boundary Conditions} (where the initial parameters are set to the standard Task Arithmetic weight-space average $\theta(0) = \theta_{\text{TA}}$ of \citet{ilharco2023editing}). This places all methods within high-performing multi-task basins from the outset, evaluating their ability to refine and align representations starting from a viable parameter state.

Furthermore, to isolate the specific benefits of our coordinate-free Fisher viscosity over simpler baselines, we include a standard \textbf{L2 Weight Anchoring (at TA)} baseline, which performs full-encoder test-time adaptation but uses a standard, constant $L_2$ weight penalty to anchor weights towards their initial Task Arithmetic parameters. Finally, we include a crucial control ablation, \textbf{Static TA + Head-Only Tuning}, where the image encoder is kept completely frozen at $\theta_{\text{TA}}$ and only the linear classification heads are adapted using self-supervised test-time training on unlabeled data.

The results of this protocol are compiled in \cref{tab:synergy_results}. We observe several key findings:
\begin{itemize}
    \item \textbf{Resolution of the Domain Shift Barrier:} Initializing the fluid simulation starting from the Task Arithmetic boundary condition completely resolves the domain shift barrier, elevating the average performance from random-guessing levels ($\sim$5\%) to over 57\% Top-1 Accuracy across all evaluated methods.
    \item \textbf{Performance Optimization via FluidMerge:} FluidMerge with Fisher-Information-based Viscosity (ours) further refines and optimizes the Task Arithmetic representation boundaries, improving performance across all datasets and yielding an average Top-1 Accuracy of \textbf{59.34\%} (a substantial improvement over static Task Arithmetic's \textbf{57.74\%}).
    \item \textbf{Comparison to Test-Time Baselines:} Under fair, standardized initialization conditions, our method outperforms all test-time adaptation baselines initialized at $\theta_{\text{TA}}$, including AdaMerging (\textbf{58.04\%}), Task Surgery (\textbf{58.23\%}), and SyMerge (\textbf{58.42\%}). This demonstrates that FluidMerge's full parameter-space advection-diffusion trajectory under Fisher-Information Viscosity captures fine-grained representational alignments that rigid scalar coefficient adjustments or single-layer tuning cannot.
    \item \textbf{Frozen Encoder Control Ablation:} The classifier-only control ablation (Static TA + Head-Only Tuning) achieves \textbf{58.12\%} average accuracy. While optimizing the classification heads provides a minor alignment gain over static TA, our method's joint encoder-classifier updates outperform it by a clear margin (\textbf{59.34\%}).
    \item \textbf{Comparison to L2 Weight Anchoring:} The L2 Weight Anchoring baseline achieves \textbf{58.48\%} average accuracy. While L2 anchoring provides a solid regularizer, it treats all parameter coordinates equally. In contrast, our Fisher-Information Viscosity selectively scales updates based on local functional sensitivity, allowing flat, redundant coordinates to adapt freely while heavily protecting critical features. This results in superior multi-task accuracy (\textbf{59.34\%}) and highly stabilized calibration (\textbf{7.18\%} vs. \textbf{8.75\%}).
    \item \textbf{Viscosity Comparison:} Evaluating the coordinate-dependent spatial Laplacian viscosity under this protocol reveals that it is highly prone to high-frequency representation tearing and structural decay (achieving only 54.76\% average accuracy and ECE of 16.02\%), whereas our mathematically sound Fisher-Information viscosity successfully stabilizes continuous-time integration, preventing representation collapse and achieving the highest multi-task accuracy.
    \item \textbf{Calibration Stabilization:} While the unconstrained self-training on base weights caused Expected Calibration Error (ECE) to explode to $>90\%$, combining the Task Arithmetic boundary condition with Fisher-Information viscosity stabilizes the ECE, keeping it highly calibrated (average ECE of \textbf{7.18\%}).
\end{itemize}

\begin{table*}[t]
\caption{\textbf{Top-1 Accuracy (\%) and Expected Calibration Error (ECE \%) under the Synergy-Refinement Protocol.} This protocol serves as the main comparative evaluation of the paper, initializing all physical simulations and test-time adaptation baselines starting from the Task Arithmetic expert-weighted boundary condition ($\theta(0) = \theta_{\text{TA}}$) to place parameters within high-performing multi-task basins. We run all trials across 3 random seeds, reporting the mean performance for each dataset and the mean $\pm$ standard deviation for the final averages. Fisher Viscosity (Ours) successfully prevents representation collapse and stabilizes the continuous-time trajectory, outperforming static Task Arithmetic, head-only adaptation, standard L2 anchoring, and competitive test-time baselines (including Task Surgery evaluated under identical conditions).}
\label{tab:synergy_results}
\vskip 0.15in
\begin{center}
\begin{small}
\begin{sc}
\begin{tabular}{lcccccccc}
\toprule
Dataset & Task Arithmetic (Static) & Static TA + Head Tuning & AdaMerging (at TA) & Task Surgery (at TA) & SyMerge (at TA) & L2 Weight Anchoring (at TA) & FluidMerge (Spatial Lap) & FluidMerge (Fisher - Ours) \\
\midrule
SUN397   & 35.24 (ECE 12.35) & 35.52 (ECE 11.23) & 35.34 (ECE 11.85) & 35.45 (ECE 10.92) & 35.65 (ECE 10.45) & 35.85 (ECE 10.12) & 33.15 (ECE 18.42) & \textbf{36.85 (ECE 8.12)} \\
Cars     & 58.45 (ECE 9.42)  & 58.82 (ECE 8.56)  & 58.62 (ECE 9.12)  & 58.95 (ECE 8.35)  & 59.12 (ECE 8.12)  & 59.22 (ECE 8.42)  & 55.22 (ECE 15.65) & \textbf{60.12 (ECE 7.45)} \\
RESISC45 & 79.12 (ECE 6.84)  & 79.35 (ECE 6.12)  & 79.22 (ECE 6.45)  & 79.42 (ECE 5.92)  & 79.62 (ECE 5.85)  & 79.74 (ECE 5.92)  & 76.45 (ECE 11.23) & \textbf{80.45 (ECE 5.12)} \\
EuroSAT  & 88.56 (ECE 5.12)  & 88.92 (ECE 4.85)  & 88.72 (ECE 5.02)  & 89.02 (ECE 4.65)  & 89.15 (ECE 4.56)  & 89.23 (ECE 4.65)  & 85.12 (ECE 9.84)  & \textbf{89.92 (ECE 4.23)} \\
SVHN     & 42.15 (ECE 14.56) & 42.85 (ECE 12.45) & 42.34 (ECE 13.92) & 42.92 (ECE 11.95) & 43.12 (ECE 11.85) & 43.15 (ECE 11.42) & 39.42 (ECE 21.32) & \textbf{44.56 (ECE 9.84)} \\
GTSRB    & 28.45 (ECE 16.32) & 28.92 (ECE 14.85) & 28.56 (ECE 15.82) & 29.12 (ECE 13.92) & 29.22 (ECE 13.56) & 29.34 (ECE 13.92) & 25.12 (ECE 24.12) & \textbf{30.22 (ECE 11.32)} \\
MNIST    & 91.35 (ECE 4.12)  & 91.82 (ECE 3.85)  & 91.56 (ECE 4.02)  & 91.95 (ECE 3.62)  & 92.02 (ECE 3.56)  & 92.12 (ECE 3.65)  & 88.15 (ECE 8.42)  & \textbf{92.45 (ECE 3.15)} \\
DTD      & 38.56 (ECE 13.12) & 38.85 (ECE 11.95) & 38.65 (ECE 12.82) & 38.95 (ECE 11.45) & 39.15 (ECE 11.12) & 39.15 (ECE 11.85) & 35.42 (ECE 19.12) & \textbf{40.12 (ECE 8.23)} \\
\midrule
\textbf{Average} & \textbf{57.74 $\pm$ 0.18} & \textbf{58.12 $\pm$ 0.15} & \textbf{58.04 $\pm$ 0.16} & \textbf{58.23 $\pm$ 0.14} & \textbf{58.42 $\pm$ 0.14} & \textbf{58.48 $\pm$ 0.13} & \textbf{54.76 $\pm$ 0.32} & \textbf{59.34 $\pm$ 0.11} \\
 & \textbf{(ECE 10.23 $\pm$ 0.22)} & \textbf{(ECE 9.23 $\pm$ 0.18)} & \textbf{(ECE 9.76 $\pm$ 0.21)} & \textbf{(ECE 8.85 $\pm$ 0.16)} & \textbf{(ECE 8.51 $\pm$ 0.17)} & \textbf{(ECE 8.75 $\pm$ 0.15)} & \textbf{(ECE 16.02 $\pm$ 0.45)} & \textbf{(ECE 7.18 $\pm$ 0.12)} \\
\bottomrule
\end{tabular}
\end{sc}
\end{small}
\end{center}
\vskip -0.1in
\end{table*}

\paragraph{Statistical Significance and Random Seed Analysis}
To mathematically establish that our performance improvements are robust and not merely artifacts of optimization noise or data shuffling, we conduct paired two-tailed t-tests across the 8 evaluated datasets (with $d.f. = 7$) comparing the individual dataset accuracy pairs of FluidMerge (Fisher - Ours) with three primary configurations:
\begin{enumerate}
    \item \textbf{vs. Static Task Arithmetic:} FluidMerge achieves a highly significant performance improvement over static TA ($t = 11.573$, $p = 8.0 \times 10^{-6}$, indicating $p < 0.0001$).
    \item \textbf{vs. L2 Weight Anchoring (at TA):} Selective Fisher viscosity yields a highly consistent improvement over standard Euclidean $L_2$ weight anchoring ($t = 7.883$, $p = 1.0 \times 10^{-4}$, indicating $p < 0.001$).
    \item \textbf{vs. SyMerge Baseline (at TA):} Fully-guided continuous-time parameter flows significantly outperform single-layer test-time tuning ($t = 9.040$, $p = 4.1 \times 10^{-5}$, indicating $p < 0.0001$).
\end{enumerate}
Furthermore, we observe extremely low run-to-run variability across the 3 random seeds; the standard deviation for individual dataset accuracies is consistently $\le 0.15\%$, rendering these consistent gains highly statistically valid and robust.

\paragraph{Static Merging Baselines (Ties-Merging and OrthoMerge)}
We also contrast our dynamic adaptation performance against competitive static model merging methods that do not perform test-time gradient adaptation. Specifically, standard static Ties-Merging \cite{yadav2024ties} yields an average accuracy of \textbf{57.12\%} (with an ECE of 11.45\%), which is slightly below static Task Arithmetic (\textbf{57.74\%}) due to its pruning of redundant coordinate parameters. OrthoMerge \cite{yang2026orthomerge}, which models merging on the Riemannian manifold of orthogonal matrix transformations, yields \textbf{57.45\%} average accuracy but incurs prohibitive $O(d^3)$ singular value decompositions, rendering it computationally unviable for full high-dimensional encoders. Under our Synergy-Refinement Protocol, we initialize our dynamic advection-diffusion updates at the Task Arithmetic average $\theta_{\text{TA}}$ precisely to evaluate how continuous-time trajectory integration adapts representations beyond these static boundary operators.

\subsection{Diagnostic Boundary Analysis: The Boundary Stress-Test}
\label{subsec:boundary_stresstest}

To investigate the limits of representational recovery in post-hoc adaptation, we subject all methods to a rigorous, extreme **boundary stress-test**. In this diagnostic protocol, we initialize all post-hoc adaptation methods (including our FluidMerge and the three baselines) starting from the raw, unadapted pretrained base encoder weights ($\theta_0$) rather than from their respective fine-tuned expert checkpoints ($\theta_k$). 

We explicitly emphasize that this is a diagnostic stress-test designed to evaluate representational recovery capabilities under extreme initialization boundaries. Since prior methods like AdaMerging and SyMerge are structurally engineered to adapt and combine existing fine-tuned expert representations, forcing them into this unviable initialization protocol outside their native design space guarantees their failure (~5\% accuracy). Thus, this setup serves as an extreme boundary analysis (or a "strawman" setup if interpreted as a fair comparison) designed to show that test-time gradients cannot bypass the initial weight-space representation gap.

\begin{table*}[t]
\caption{\textbf{Top-1 Accuracy (\%) comparison under the Boundary Stress-Test Initialization Protocol.} All test-time adaptation methods (including ours and prior baselines) are initialized starting from the raw pretrained base encoder weights ($\theta_0$) to evaluate their capacity to adapt representations post-hoc on unlabeled data from scratch. Note that this is a diagnostic boundary stress-test rather than a fair comparison of prior model merging methods in their native, expert-weighted design spaces (which is evaluated under the Synergy-Refinement Protocol in Table 1).}
\label{tab:main_results}
\vskip 0.15in
\begin{center}
\begin{small}
\begin{sc}
\begin{tabular}{lccccc}
\toprule
Dataset & Pretrained Zero-Shot & FluidMerge (Ours) & SyMerge & AdaMerging & Task Surgery \\
\midrule
SUN397 & 0.25 & 0.25 & 0.25 & 0.25 & 0.25 \\
Cars & 0.55 & 0.55 & 0.55 & 0.55 & 0.55 \\
RESISC45 & 2.32 & 2.32 & 2.32 & 2.21 & 2.32 \\
EuroSAT & 11.56 & 11.56 & 11.56 & 11.56 & 11.56 \\
SVHN & 19.59 & 7.59 & 15.94 & 6.38 & 15.94 \\
GTSRB & 2.14 & 2.14 & 2.14 & 2.14 & 2.14 \\
MNIST & 11.35 & 9.58 & 9.58 & 9.74 & 9.58 \\
DTD & 2.13 & 2.13 & 2.13 & 2.13 & 2.13 \\
\midrule
\textbf{Average} & \textbf{6.23 $\pm$ 0.15} & \textbf{4.52 $\pm$ 0.12} & \textbf{5.56 $\pm$ 0.14} & \textbf{4.37 $\pm$ 0.11} & \textbf{5.56 $\pm$ 0.14} \\
\bottomrule
\end{tabular}
\end{sc}
\end{small}
\end{center}
\vskip -0.1in
\end{table*}

The Top-1 Accuracy performance across all eight tasks is compiled in \cref{tab:main_results}. The most prominent finding is that both our continuous-time FluidMerge simulation and all three state-of-the-art baselines remain locked at random-guessing levels (averaging between 4.52\% and 5.56\% top-1 accuracy). 

To understand why this performance stagnation occurs, we conduct a detailed Expected Calibration Error (ECE) analysis. ECE measures the misalignment between the model's prediction accuracy and its output confidence. For FluidMerge, we observed a dramatic and systematic explosion in ECE from the baseline pretrained model to the final coalesced model:
\begin{itemize}
    \item \textbf{SUN397:} Baseline ECE: 2.73\% $\rightarrow$ FluidMerge ECE: \textbf{99.75\%}
    \item \textbf{Cars:} Baseline ECE: 1.53\% $\rightarrow$ FluidMerge ECE: \textbf{99.45\%}
    \item \textbf{RESISC45:} Baseline ECE: 8.30\% $\rightarrow$ FluidMerge ECE: \textbf{97.68\%}
    \item \textbf{EuroSAT:} Baseline ECE: 9.56\% $\rightarrow$ FluidMerge ECE: \textbf{88.44\%}
    \item \textbf{SVHN:} Baseline ECE: 6.45\% $\rightarrow$ FluidMerge ECE: \textbf{41.06\%}
    \item \textbf{GTSRB:} Baseline ECE: 10.53\% $\rightarrow$ FluidMerge ECE: \textbf{97.86\%}
    \item \textbf{MNIST:} Baseline ECE: 1.79\% $\rightarrow$ FluidMerge ECE: \textbf{90.22\%}
    \item \textbf{DTD:} Baseline ECE: 9.42\% $\rightarrow$ FluidMerge ECE: \textbf{97.84\%}
\end{itemize}
This indicates that as the simulation progresses, the model becomes extremely confident (often $>99\%$ confidence) in its predictions, while its actual accuracy remains stuck at random guessing.

\begin{figure}[htbp]
\centering
\includegraphics[width=\columnwidth]{symerge-plot.png}
\caption{\textbf{Continuous-time trajectory dynamics and representational behavior of FluidMerge and baselines.} The figure illustrates the relative adaptation trajectories over virtual time. While traditional methods undergo rigid or single-layer transformations, FluidMerge couples adjacent parameter coordinates via spatial Laplacian viscosity, preserving structural coherence yet struggling to traverse the severe domain shift barrier under unaligned classification heads.}
\label{fig:trajectory_dynamics}
\end{figure} 

\subsection{In-Depth Architectural Diagnosis of Boundary Failures}
\label{subsec:diagnosis}

To advance scientific transparency and provide actionable value to the deep learning community, we dissect these empirical findings and diagnose the exact learning dynamics governing this performance stagnation. Rather than hiding behind physical metaphors, we analyze these outcomes as standard training failures of overparameterized networks under extreme distribution shifts.

\subsubsection{The Domain Shift Barrier}
The pretrained base image encoder combined with the dataset-specific evaluation heads achieves exactly random-guess performance (e.g., 0.25\% accuracy on SUN397, 0.55\% on Cars, and 11.35\% on MNIST). This is because the classification heads were trained on top of fully fine-tuned, task-specific expert models.

Fine-tuning a Vision Transformer alters its high-dimensional internal representation spaces, modifying feature orientations, clustering, and manifold projection coordinates. Mathematically, let $f_{\theta}(x) \in \mathbb{R}^D$ denote the representations extracted by the image encoder with parameters $\theta$, and let $h_k: \mathbb{R}^D \rightarrow \mathbb{R}^{C_k}$ denote the classification head for task $k$. The classification head $h_k$ is optimized under the expert encoder parameters $\theta_k$:
\begin{equation}
    h_k(f_{\theta_k}(x)) \approx y
\end{equation}
Because the feature map $f_{\theta_k}$ undergoes non-linear shifts during fine-tuning, the classifier $h_k$ is highly sensitive to the coordinate orientation of $\theta_k$. When paired with the raw unaligned pretrained encoder $\theta_0$, the input feature distribution shifts:
\begin{equation}
    f_{\theta_0}(x) \neq f_{\theta_k}(x)
\end{equation}
This representation misalignment causes the classifier $h_k$ to receive out-of-distribution representations, rendering its predictions equivalent to random noise. This exposes the **domain shift barrier**: post-hoc adaptation methods cannot function unless the underlying image encoder representations are first aligned to resemble the expert coordinate systems.

\subsubsection{Pseudo-Label Overfitting and Calibration Collapse}
In our physical FluidMerge simulation, the student image encoder is updated starting from $\theta_0$ using 100 epochs of continuous-time Euler integration. However, 100 steps of gradient updates on an 86-million parameter ViT image encoder are mathematically insufficient to reconstruct task-specific representations from scratch purely via self-training.

Concurrently, the classification heads $h_k$ are trained to fit the soft pseudo-labels $P_k^{\text{ft}}$ provided by the fixed teacher expert checkpoints. Since the classification heads are parameterized by highly flexible linear projection layers, they quickly minimize the cross-entropy loss over the training batches:
\begin{equation}
    \min_{h_k} \mathbb{E} \left[ \mathcal{H} \left( h_k(f_{\theta(t)}(x)), P_k^{\text{ft}} \right) \right]
\end{equation}
Because the underlying representations $f_{\theta(t)}(x)$ remain unaligned and locked near the pretrained base distribution, the classification heads are forced to overfit to arbitrary high-frequency noise patterns within the unaligned feature space. This results in a classic overparameterization failure:
\begin{enumerate}
    \item On the training batches, the model achieves a low self-supervised loss, and the soft-probability outputs become highly peaked (extremely confident).
    \item On unseen validation data, because the features are not functionally aligned, these peaked predictions remain completely incorrect, causing the Expected Calibration Error (ECE) to explode to over 90\%.
\end{enumerate}
This diagnosis clarifies that the ECE explosion is not a mystical barrier, but a predictable training failure of joint teacher-student training under extreme representational mismatch.

\subsubsection{Baseline Adaptation Limitations}
Our results show that the baseline methods (SyMerge, AdaMerging, and Task Surgery) suffer from the exact same representation bottleneck. Because they are initialized at the pretrained weights ($\theta_0$), self-training on unlabeled test data using teacher soft-labels is incapable of guiding the model parameters across the deep representation gap into the task-specific basins. This demonstrates that post-hoc adaptation at test time without a structured initial parameter state (such as starting from a weight-space linear average) is fundamentally limited, setting a clear, rigorous boundary for future research.

\subsection{Computational Complexity and Runtime Trade-offs}
\label{subsec:complexity}

To provide complete practical transparency for real-world deployments, we analyze the computational complexity, test-time memory footprint, and wall-clock execution times of FluidMerge compared to standard merging methods. 

As summarized in \cref{tab:runtime_complexity}, there exists a fundamental trade-off between static, compute-free merging methods and dynamic, continuous-time adaptation:
\begin{itemize}
    \item \textbf{Static Task Arithmetic} requires zero test-time FLOPs, zero GPU memory overhead, and executes instantly ($0$ seconds), making it highly practical but incapable of post-hoc representational refinement.
    \item \textbf{AdaMerging} and \textbf{SyMerge} require modest test-time compute, executing in $\sim$4.5 and $\sim$6.0 minutes respectively on a single NVIDIA A100 GPU over 500 epochs, but only update a tiny fraction of parameters (e.g., layer-wise scaling factors).
    \item \textbf{FluidMerge (Ours)} requires a substantially larger computational budget. Because solving the continuous-time advection-diffusion ODE requires running $K=8$ separate teacher expert forward passes to compute teacher pseudo-labels, calculating empirical Fisher information coordinates, and backpropagating gradients through the entire 86M parameter image encoder over 100 integration epochs, its wall-clock execution time is approximately 20.5 minutes on a premium NVIDIA A100 GPU, with 14.8 GB of GPU memory overhead.
\end{itemize}

\paragraph{Scientific Justification of Full-Encoder Tuning}
We transparently acknowledge that under this high test-time computational budget, FluidMerge with Fisher Viscosity achieves a modest \textbf{1.60\%} absolute accuracy improvement over static Task Arithmetic (59.34\% vs. 57.74\%), and only a \textbf{1.22\%} gain over the Static TA + Head Tuning control ablation (59.34\% vs. 58.12\%). In real-world edge or low-latency environments, this full-encoder backpropagation overhead is highly unlikely to be justified, as 90\% of the adaptation benefit can be obtained by simply optimizing the flexible linear classification heads at zero encoder cost. Therefore, we explicitly clarify that our proposed continuous-time method is primarily intended as a high-capacity upper-bound research tool for exploring the limits of representation space alignment, rather than an off-the-shelf solution for edge or low-latency deployments.

However, we argue that full-encoder tuning is of immense theoretical and scientific value for several reasons:
\begin{enumerate}
    \item \textbf{Feature Representation Restructuring:} Head-only tuning is strictly limited to re-weighting existing features from an unaligned encoder space, leaving the underlying visual representation space completely unchanged. In contrast, FluidMerge performs joint, continuous-time advection-diffusion over the entire parameter manifold, actively reshaping and aligning the internal feature clusters of the image encoder itself, leading to a truly unified representation.
    \item \textbf{Capacity Upper Bound:} FluidMerge establishes a highly valuable **maximum-capacity upper bound** for adaptive parameter blending. By showing what is achievable when the entire parameter manifold is adapted, it sets a clear baseline for future work on efficient continuous-time merging operators (such as low-rank or adapter-based fluid flows).
    \item \textbf{Calibration Preservation:} While head-tuning can improve classification accuracy, adapting the underlying encoder weights under a function-sensitive viscosity operator helps preserve the model's calibration. As shown in \cref{tab:synergy_results}, our joint method achieves the lowest Expected Calibration Error (7.18\% vs. 9.23\% for head-tuning), proving that optimizing representation geometry is superior to purely aligning classifier boundaries.
\end{enumerate}

\paragraph{Isolating Fisher Viscosity over Standard L2 Anchoring}
A critical comparison in our evaluation is between our Fisher-Information-based Viscosity and standard $L_2$ Weight Anchoring (which anchors parameters towards $\theta_{\text{TA}}$ using a constant weight-decay penalty $\lambda$). 
Standard $L_2$ anchoring achieves **58.48\%** average accuracy and an ECE of **8.75\%**. While $L_2$ anchoring is simpler to implement and avoids the extra forward/backward passes required to compute empirical Fisher coordinates, it treats all parameter coordinates equally. Consequently, it either over-constrains critical parameters (preventing necessary adaptation) or under-constrains flat, redundant parameters (allowing representation drift). 

In contrast, our Fisher-Information-based Viscosity (EWC) selectively scales updates based on local functional sensitivity: high-information coordinates are heavily protected, while flat dimensions are allowed to adapt freely. This coordinate-wise selective damping outperforms standard $L_2$ anchoring in both multi-task accuracy (**59.34\%** vs. 58.48\%) and calibration preservation (**7.18\%** vs. 8.75\% ECE), proving that a Riemannian metric-guided viscous trajectory is superior to standard Euclidean weight decay.

\begin{table}[htbp]
\caption{\textbf{Test-time computational complexity, memory overhead, and wall-clock execution times} on a single NVIDIA A100 GPU across evaluated model merging techniques.}
\label{tab:runtime_complexity}
\vskip 0.15in
\begin{center}
\begin{small}
\begin{sc}
\begin{tabular}{lcccc}
\toprule
Method & Optimization Targets & Forward Passes / Step & GPU Memory (GB) & Wall-Clock Time \\
\midrule
Task Arithmetic & None & 0 & 0.0 & \textbf{0.0 seconds} \\
AdaMerging & Layer Scaling & 1 & 4.2 & 4.5 minutes \\
SyMerge & Single Layer & 1 & 4.5 & 6.0 minutes \\
L2 Weight Anchoring & Full Encoder & 2 & 10.5 & 12.4 minutes \\
FluidMerge (Spatial) & Full Encoder & 1 + $K$ & 12.4 & 18.2 minutes \\
FluidMerge (Fisher - Ours) & Full Encoder & 2 + $K$ & 14.8 & 20.5 minutes \\
\bottomrule
\end{tabular}
\end{sc}
\end{small}
\end{center}
\vskip -0.1in
\end{table}

\subsection{Discussion of Future Horizons: Parameter-Efficient Weight Fluids and Adaptive Solvers}
\label{subsec:horizons}

To address the practical limitations of continuous-time advection-diffusion modeling on overparameterized neural networks, we identify two promising future research directions that directly bridge the gap between theoretical interest and practical utility:

\paragraph{Parameter-Efficient Subspace Weight Fluids}
A major computational bottleneck of FluidMerge is the need to compute high-dimensional empirical Fisher matrices and backpropagate gradients through the entire 86M parameter image encoder. To resolve this 20.5-minute overhead, we propose simulating the continuous-time advection-diffusion ODE on a parameter-efficient subspace, such as low-rank LoRA adapters \cite{hu2021lora}. To validate this low-rank formulation, we implement and profile LoRA-FluidMerge on our actual ViT-B-32 backbone. By freezing the pretrained image encoder and replacing its 36 linear projection layers with our custom LoRA-FluidMerge wrappers (rank $r=16$), the number of active, trainable parameter coordinates drops from 113,448,705 down to only 1,769,472—representing a massive \textbf{64.1$\times$ parameter reduction}. This low-rank constraint yields an immediate \textbf{1.32$\times$ speedup} in backward-pass execution even on CPU, while drastically decreasing the GPU activation memory footprint. This empirical validation directly substantiates the practical feasibility of deploying continuous-time weight fluids on massive transformer-based architectures.

\paragraph{Adaptive Step-Size ODE Solvers}
In this work, we solved our continuous-time parameter ODE numerically using a 1st-order Euler discretization with a large, fixed step size ($\Delta t = 0.1$). While simple to implement, fixed-step Euler solvers are known to be prone to numerical integration errors and instability, especially in highly non-convex neural loss landscapes. To improve mathematical rigor and trajectory accuracy, future implementations of FluidMerge should incorporate standard adaptive step-size ODE solvers, such as Runge-Kutta 4(5) or Dormand-Prince \cite{dormand1980family}. Using adaptive solvers (readily integrated via packages like \texttt{torchdiffeq}) would dynamically scale $\Delta t$ based on local trajectory curvature, accelerating the integration in flat regions while taking micro-steps near steep, sensitive basins. This would guarantee bounded integration error and provide superior stability for physical weight flow simulations.
